%==============================================================================
%  paper.tex — Tesis (formato IEEE) en español
%
%  Compila con:
%      pdflatex paper.tex
%      pdflatex paper.tex      (segunda pasada para resolver referencias)
%
%  Las figuras están en:
%      ../../data/derived/figures/section25/
%  (rutas relativas desde docs/paper/paper.tex)
%
%  Si LaTeX no encuentra IEEEtran.cls instala el paquete texlive-publishers
%  o coloca IEEEtran.cls en el mismo directorio que este archivo.
%==============================================================================
\documentclass[journal,onecolumn,11pt]{IEEEtran}

% --- Codificación y idioma ---
\usepackage[utf8]{inputenc}
\usepackage[T1]{fontenc}
\usepackage[spanish,es-tabla,es-noindentfirst]{babel}
\usepackage{lmodern}

% --- Matemáticas ---
\usepackage{amsmath,amssymb,amsfonts,mathtools}
\usepackage{bm}

% --- Gráficos y tablas ---
\usepackage{graphicx}
\usepackage{booktabs}
\usepackage{multirow}
\usepackage{array}
\usepackage{siunitx}
\sisetup{group-separator={\,},output-decimal-marker={,},detect-weight=true}

% --- Diagramas de bloque ---
\usepackage{tikz}
\usetikzlibrary{positioning,arrows.meta,calc,shapes.geometric}

% --- Hipervínculos ---
\usepackage[hidelinks]{hyperref}

% --- Pseudocódigo ---
\usepackage{algorithm}
\usepackage{algpseudocode}
\floatname{algorithm}{Algoritmo}
\renewcommand{\algorithmicrequire}{\textbf{Entrada:}}
\renewcommand{\algorithmicensure}{\textbf{Salida:}}
\renewcommand{\algorithmicwhile}{\textbf{mientras}}
\renewcommand{\algorithmicdo}{\textbf{hacer}}
\renewcommand{\algorithmicend}{\textbf{fin}}
\renewcommand{\algorithmicfor}{\textbf{para}}
\renewcommand{\algorithmicif}{\textbf{si}}
\renewcommand{\algorithmicthen}{\textbf{entonces}}
\renewcommand{\algorithmicelse}{\textbf{si no}}
\renewcommand{\algorithmicreturn}{\textbf{devolver}}

% --- Macros útiles ---
\newcommand{\R}{\mathbb{R}}
\newcommand{\E}{\mathbb{E}}
\newcommand{\xv}{\mathbf{x}}
\newcommand{\zv}{\mathbf{z}}
\newcommand{\Fm}{\mathbf{F}}
\newcommand{\Hm}{\mathbf{H}}
\newcommand{\Qm}{\mathbf{Q}}
\newcommand{\Rm}{\mathbf{R}}
\newcommand{\Pm}{\mathbf{P}}
\newcommand{\Km}{\mathbf{K}}
\newcommand{\Iv}{\mathbf{I}}

\graphicspath{{../../data/derived/figures/section25/}}

%==============================================================================
\begin{document}

\title{Ejecución adaptativa en mercados financieros con filtro de Kalman y\\
       algoritmo Hedge: una arquitectura de control para BTC/USDT}

\author{Damián Muñoz Díaz%
\thanks{Manuscrito preparado en mayo de 2026.
        Repositorio del proyecto: \texttt{adversity-is-controllable}.}}

\markboth{Tesis --- Ingeniería de Control y Microestructura de Mercado}{}

\maketitle

%==============================================================================
\begin{abstract}
Este trabajo describe una arquitectura adaptativa para ejecutar órdenes de
compra de Bitcoin contra Tether (par BTC/USDT del exchange Binance) usando
datos reales de mercado a una resolución de 100 ms. El sistema combina un
filtro de Kalman lineal --- que estima dos variables ocultas del mercado:
la presión direccional y el régimen de volatilidad --- con un algoritmo
de aprendizaje en línea conocido como Hedge, que selecciona en cada tick
una de tres acciones: esperar, postear una orden límite, o cruzar el spread
con una orden a mercado. Sobre cuatro sesiones independientes que totalizan
217\,425 ticks de mercado, se evaluaron cuatro variantes del controlador:
una versión sin estado, una con condicionamiento unidimensional sobre la
presión, una bidimensional con presión más régimen, y tres bidimensionales
adversariales que reemplazan el régimen por señales rolling de 60 segundos.
El resultado central es que la versión unidimensional condicionada en presión
gana en todos los pares estadísticamente significativos (160 corridas
comparadas, $|t|\!\geq\!17$ frente a cualquier competidor multidimensional)
y que la causa mecanística del fracaso de las variantes ricas en estado es
una \emph{barrera de toxicidad}: cualquier eje secundario que detecta
oportunidades de fill aumenta la fracción de órdenes pasivas, las cuales
se ejecutan precisamente en los instantes previos a movimientos
desfavorables del precio. Se proponen las direcciones siguientes para la
tesis: (i) condicionamiento por tamaño de orden y (ii) confirmación
multi-sesión del resultado positivo de presión.
\end{abstract}

\begin{IEEEkeywords}
Microestructura de mercado, filtro de Kalman, aprendizaje en línea,
algoritmo Hedge, ejecución óptima, BTC/USDT, criptomonedas, control
adaptativo, libro de órdenes.
\end{IEEEkeywords}

%==============================================================================
\section{Introducción}\label{sec:intro}

Cualquier persona que compre una criptomoneda en una bolsa electrónica
enfrenta un problema simple de enunciar pero difícil de resolver:
\emph{cómo} comprar. Si se cruza el diferencial de inmediato (\emph{orden
a mercado}) la compra es segura pero costosa porque se paga el lado más
caro del libro. Si se postea una orden límite al mejor bid el precio
medio queda mejor pero la orden puede no ejecutarse; y cuando se ejecuta
suele ser justo antes de que el precio caiga --- es decir, justo cuando
peor compra hace el comprador. Si se espera, se corre el riesgo de que
el precio suba y se compre más caro un instante después.

\medskip

La tesis aborda este problema bajo el supuesto operativo siguiente: en
cada uno de los más de \num{217000} ticks que componen las cuatro
sesiones experimentales, un agente está \emph{obligado} a comprar 1 BTC.
La obligación elimina del problema la decisión de cuánto comprar y deja
solo la decisión de \emph{cómo} hacerlo. Las tres acciones son
\textsc{Wait}, \textsc{Passive} (orden límite al best bid) y
\textsc{Aggressive} (orden a mercado al best ask).

\medskip

La pregunta de investigación, en una sola línea, es:
\begin{quote}
\emph{Dadas las condiciones actuales del mercado, ¿cómo debe actuar el
agente para minimizar el costo conjunto de slippage y movimiento
adverso?}
\end{quote}

\medskip

La hipótesis central es que las condiciones actuales del mercado pueden
resumirse en dos números ocultos --- la \emph{presión} (cuánto empuje
direccional hay) y el \emph{régimen} (cuánta turbulencia hay) --- y que
un filtro de Kalman puede estimarlos a partir de tres observaciones
extraídas del libro de órdenes. Una vez estimados, un algoritmo de
aprendizaje en línea (Hedge) elige la acción que minimiza el costo
acumulado de ejecución.

\medskip

El resto del artículo se organiza así. La Sección~\ref{sec:mercado}
explica desde cero la mecánica de compra en BTC/USDT y todos los
términos económicos que aparecen después. La Sección~\ref{sec:arq}
describe la arquitectura del sistema con un diagrama de bloques.
La Sección~\ref{sec:features} desglosa las features observables.
La Sección~\ref{sec:kalman} \emph{deconstruye número por número} las
matrices del filtro de Kalman. La Sección~\ref{sec:hedge} explica el
algoritmo Hedge y sus tres variantes condicionadas por estado.
La Sección~\ref{sec:simulador} describe el simulador de ejecución y
la función de pérdida. La Sección~\ref{sec:experimentos} presenta el
diseño experimental y la Sección~\ref{sec:resultados} los resultados
con las cinco figuras principales. La Sección~\ref{sec:discusion}
explica el hallazgo mecanístico central, llamado \emph{barrera de
toxicidad}, y la Sección~\ref{sec:conclusiones} concluye y sugiere
los próximos pasos.

%==============================================================================
\section{Mercado y mecánica de ejecución en BTC/USDT}\label{sec:mercado}

Esta sección está escrita pensando en una persona con formación de
ingeniería pero sin experiencia previa en mercados financieros. El
objetivo es que cuando aparezca \emph{spread} o \emph{OFI} más
adelante en el texto, el lector ya sepa exactamente qué significa.

\subsection{El par BTC/USDT}
BTC es el ticker de Bitcoin. USDT es Tether, una criptomoneda diseñada
para mantener un valor cercano a un dólar estadounidense. El par
\emph{BTC/USDT} indica el precio de un Bitcoin medido en USDT.
Una cotización del estilo
\[
\text{BTC/USDT} = 76\,230{,}70
\]
significa que en el momento de la cotización un Bitcoin se cambia por
unos \num{76230,70} USDT. Como un USDT es aproximadamente un dólar,
esto se interpreta también como ``el Bitcoin vale unos 76\,230 dólares''.

\subsection{El libro de órdenes (\emph{order book})}

Una bolsa electrónica como Binance funciona con un \emph{libro de
órdenes} que es esencialmente dos listas:

\begin{itemize}
\item Los \emph{bids} son las órdenes de compra que están esperando.
      Cada bid tiene un \textbf{precio} (cuánto está dispuesto a pagar
      el comprador por 1 BTC) y una \textbf{cantidad} (cuántos BTC
      quiere a ese precio). Los bids se ordenan de mayor a menor precio:
      el bid con el precio más alto es el \emph{best bid}.
\item Los \emph{asks} (también llamados \emph{offers}) son las órdenes
      de venta. Tienen también precio y cantidad, ordenadas de menor a
      mayor precio: el ask con precio más bajo es el \emph{best ask}.
\end{itemize}

\noindent
Una representación cualitativa del libro a un instante $t$ se muestra
en la Tabla~\ref{tab:libro-ejemplo}.

\begin{table}[!htbp]
\caption{Ejemplo simplificado del libro de órdenes.}
\label{tab:libro-ejemplo}
\centering
\begin{tabular}{rrcll}
\toprule
\multicolumn{2}{c}{\textbf{Bids (compra)}} & & \multicolumn{2}{c}{\textbf{Asks (venta)}} \\
\cmidrule(r){1-2} \cmidrule(l){4-5}
Precio (USDT) & Cantidad (BTC) & & Precio (USDT) & Cantidad (BTC) \\
\midrule
76\,230{,}69 & 0{,}210 & \textsc{best} & 76\,230{,}70 & 0{,}085 \\
76\,230{,}68 & 1{,}040 &              & 76\,230{,}71 & 0{,}450 \\
76\,230{,}67 & 0{,}330 &              & 76\,230{,}72 & 0{,}120 \\
$\vdots$     & $\vdots$ &              & $\vdots$     & $\vdots$ \\
\bottomrule
\end{tabular}
\end{table}

El número clave que se deriva del libro es el \textbf{precio medio}
o \emph{mid price}:
\begin{equation}
\text{mid}_t = \frac{\text{best\_bid}_t + \text{best\_ask}_t}{2}.
\label{eq:mid}
\end{equation}

En BTC/USDT el incremento mínimo de precio (\emph{tick size}) es de
0{,}01 USDT, así que el spread más estrecho posible es de un centavo.
La realidad observada en \num{38973} ticks limpios entre el 20 y 21
de abril de 2026 es que el spread está exactamente en 0{,}01 USDT
durante el 99{,}5\% del tiempo, con picos ocasionales de hasta 8 USDT
durante eventos de baja liquidez.

\subsection{Las tres acciones del agente}

El agente, en cada tick, escoge una de tres acciones:

\begin{description}
\item[\textsc{Wait}.] No envía ninguna orden. Espera al siguiente tick
      (típicamente 100 ms después). El costo de esperar es que si el
      precio sube, en el próximo tick se compra más caro.
\item[\textsc{Passive}.] Postea una orden límite de compra de 1 BTC
      al precio del best bid actual. Esa orden \emph{solo} se ejecuta
      si en el siguiente tick aparece un vendedor que crea una orden
      de venta a ese precio (o por debajo) que cruce con la nuestra.
      En la práctica, esto ocurre cuando el precio medio cae. Si el
      precio sube, la orden no se ejecuta.
\item[\textsc{Aggressive}.] Envía una orden a mercado de 1 BTC. La
      orden se ejecuta inmediatamente contra el best ask. La compra
      siempre se concreta. Pero pagamos el lado más caro del libro:
      pagamos por encima del mid en una cantidad igual a medio spread.
\end{description}

\subsection{Slippage y adverse move (definiciones formales)}

Las dos componentes del costo de ejecución son:

\paragraph{Slippage.}
Es la diferencia entre el precio al que ejecutamos y el precio medio
en el momento de la decisión. Si la orden se ejecutó con precio
$p^{\text{fill}}$ y en el momento de decidir el mid era
$\text{mid}_t$:
\begin{equation}
\text{slippage}_t \;=\; p^{\text{fill}}_t - \text{mid}_t.
\label{eq:slippage}
\end{equation}
Para una orden agresiva, $p^{\text{fill}} = \text{best\_ask}$ y el
slippage es positivo (pagamos más que el mid). Para una orden pasiva
ejecutada, $p^{\text{fill}} = \text{best\_bid}$ y el slippage es
\emph{negativo} (compramos por debajo del mid: capturamos medio
spread). Si la orden no se ejecuta el slippage es cero por convención.

\paragraph{Adverse move.}
Mide cuánto se movió el precio en contra de nuestra posición
\emph{después} de que actuamos. Sea
$\text{next\_mid}_t = \text{mid}_{t+1}$ el precio medio en el tick
siguiente.

\begin{itemize}
\item Si la orden se ejecutó comprando a precio $p^{\text{fill}}_t$:
\begin{equation}
\text{adverse}_t = \max\bigl(0,\; p^{\text{fill}}_t - \text{next\_mid}_t\bigr).
\label{eq:adverse-fill}
\end{equation}
Es decir, si el precio cayó por debajo de lo que pagamos, registramos
la diferencia. Si el precio subió, el adverse es cero.

\item Si la orden no se ejecutó (acción \textsc{wait} o \textsc{passive}
sin fill):
\begin{equation}
\text{adverse}_t = \max\bigl(0,\; \text{next\_mid}_t - \text{mid}_t\bigr).
\label{eq:adverse-nofill}
\end{equation}
Si el precio subió, perdimos una oportunidad de comprar más barato.
Si bajó, esperar fue lo correcto y el adverse es cero.
\end{itemize}

\paragraph{Función de pérdida del tick.}
La pérdida que entra al algoritmo Hedge en cada tick es:
\begin{equation}
L_t \;=\; \text{slippage}_t \;+\; \lambda \cdot \text{adverse}_t,
\label{eq:loss}
\end{equation}
con $\lambda = 0{,}1$ por defecto. La elección de
$\lambda$ se discute en la Sección~\ref{sec:simulador}.

\subsection{Order Flow Imbalance (OFI)}\label{ssec:ofi}

OFI mide el cambio neto de presión en el libro entre dos snapshots
consecutivos. Para el nivel $i$ (donde $i\!=\!0$ es el top of book):
\begin{equation}
\mathrm{OFI}_{t,i} \;=\; \bigl(q^{\text{bid}}_{t,i} - q^{\text{bid}}_{t-1,i}\bigr)
                       - \bigl(q^{\text{ask}}_{t,i} - q^{\text{ask}}_{t-1,i}\bigr),
\label{eq:ofi}
\end{equation}
donde $q^{\text{bid}}_{t,i}$ es la cantidad agregada en el bid de
rango $i$ en el tick $t$. Si OFI$_{t,1}>0$ significa que entre
$t\!-\!1$ y $t$ creció la cantidad demandada al best bid o decreció
la ofrecida al best ask: hay presión compradora. Si OFI$<0$ ocurre
lo contrario: presión vendedora. La feature \texttt{ofi\_l1} usada en
la observación del filtro corresponde a $i\!=\!1$ (top of book). Esta
formulación sigue la convención multinivel de Bouchaud et~al.~\cite{bouchaud2018trades}.

\subsection{Depth imbalance}\label{ssec:dimbalance}

Depth imbalance agrega cantidades a través de los $N$ niveles
superiores del libro:
\begin{equation}
\text{depth\_imb}_t \;=\; \frac{Q^{\text{bid}}_t - Q^{\text{ask}}_t}
                                    {Q^{\text{bid}}_t + Q^{\text{ask}}_t}
                                    \in (-1, +1),
\label{eq:depthimb}
\end{equation}
con
$Q^{\text{bid}}_t = \sum_{i=1}^{N}\,w_i \, q^{\text{bid}}_{t,i}$ y
análogamente $Q^{\text{ask}}_t$. Los pesos lineales $w_i = N\!-\!i\!+\!1$
le dan más peso al top of book y menos a los niveles profundos.
Por defecto $N=10$. Los valores cercanos a $+1$ indican mucho más bid
que ask (presión compradora persistente); valores cercanos a $-1$
indican lo contrario.

\subsection{Volatilidad realizada vol\_30s}\label{ssec:vol30s}

\texttt{vol\_30s} es la desviación estándar de los retornos
logarítmicos del mid price calculados sobre los últimos 30 segundos:
\begin{equation}
\text{vol\_30s}_t \;=\; \mathrm{std}\!\left\{
\ln \frac{\text{mid}_s}{\text{mid}_{s-1}}
\;:\; s\in (t-30\text{s},\,t]
\right\}.
\label{eq:vol30s}
\end{equation}
Se eligió 30~s en lugar de la ventana original de 5~s porque
\texttt{vol\_5s} estaba en cero el 54\% del tiempo (el mid no se mueve
durante muchos intervalos cortos en BTC) mientras que
\texttt{vol\_30s} es no-cero el 91{,}47\% del tiempo. Más densidad de
señal $\Rightarrow$ mejor condicionamiento para el filtro.

\subsection{Microprecio}

Como referencia auxiliar usamos también el \emph{microprecio}, que
es el mid ponderado por cantidades:
\begin{equation}
\text{microprice}_t \;=\;
\frac{q^{\text{bid}}_{t,1}\,\text{best\_ask}_t \;+\;
      q^{\text{ask}}_{t,1}\,\text{best\_bid}_t}
     {q^{\text{bid}}_{t,1} + q^{\text{ask}}_{t,1}}.
\label{eq:microprice}
\end{equation}
Si en el best bid hay 0{,}1 BTC y en el best ask hay 10 BTC, el
microprecio se desplaza hacia el bid: el lado delgado predice la
dirección probable del próximo movimiento. No se usa como observación
del filtro pero se calcula y se almacena para diagnóstico.

%==============================================================================
\section{Arquitectura del sistema}\label{sec:arq}

El sistema es una tubería (\emph{pipeline}) de seis pasos
diferenciados, ilustrada en la Fig.~\ref{fig:bloques}.

\begin{figure*}[!htbp]
\centering
\begin{tikzpicture}[
  node distance=0.55cm and 0.85cm,
  block/.style={rectangle,draw,thick,rounded corners,
                minimum width=2.6cm,minimum height=1cm,align=center,
                font=\small},
  data/.style ={rectangle,draw,thick,fill=gray!12,
                minimum width=2.4cm,minimum height=0.85cm,align=center,
                font=\small\itshape},
  arrow/.style={-Stealth,thick}
  ]

  \node[data]                                       (raw)   {Binance WS\\depth + trades};
  \node[block,right=of raw]                          (book)  {Paso 1\,--\,2\\Reconstrucción\\del libro};
  \node[block,right=of book]                         (feat)  {Paso 3\\Feature pipeline};
  \node[data,right=of feat]                          (z)     {$\zv_t \in \R^3$};
  \node[block,below=of z]                            (kal)   {Paso 4\\Filtro de Kalman};
  \node[data,left=of kal]                            (xv)    {$\xv_t = (\text{pres},\text{reg})$};
  \node[block,left=of xv]                            (pol)   {Paso 5\\Política Hedge};
  \node[data,below=of pol]                           (act)   {Acción\,$\in\{\textsc{Wait},\textsc{Pass},\textsc{Agg}\}$};
  \node[block,right=of act]                          (sim)   {Paso 6\\Simulador};
  \node[data,right=of sim]                           (loss)  {Loss $L_t$};

  \draw[arrow] (raw) -- (book);
  \draw[arrow] (book) -- (feat);
  \draw[arrow] (feat) -- (z);
  \draw[arrow] (z) -- (kal);
  \draw[arrow] (kal) -- (xv);
  \draw[arrow] (xv) -- (pol);
  \draw[arrow] (pol) -- (act);
  \draw[arrow] (act) -- (sim);
  \draw[arrow] (sim) -- (loss);
  \draw[arrow] (loss) |- ++(-13.5,-0.7) |- (pol.south);

\end{tikzpicture}
\caption{Diagrama de bloques del sistema. La salida del simulador
        (loss escalar) se realimenta al algoritmo Hedge para actualizar
        las probabilidades de acción.}
\label{fig:bloques}
\end{figure*}

\medskip

A continuación se describe cada bloque en una línea para que el
lector tenga el mapa completo antes de la matemática.

\begin{enumerate}
\item \textbf{Pasos 1--2 --- Ingesta y libro.} Un cliente WebSocket lee
      el flujo de actualizaciones de Binance y reconstruye el libro
      local mediante un \texttt{SortedDict} con claves negadas
      (los bids como \texttt{-precio} para que el índice 0 sea siempre
      el mejor bid). Periódicamente se sincroniza con un snapshot REST
      para defenderse de cortes de WebSocket; ante una brecha de
      identificadores de actualización se hace \emph{self-heal} desde
      el snapshot que la cubre.
\item \textbf{Paso 3 --- Feature pipeline.} Cada snapshot del libro
      pasa por funciones puras que calculan
      \texttt{depth\_imbalance}, \texttt{ofi\_l1},
      \texttt{spread\_abs}, \texttt{vol\_30s} y \texttt{microprice}.
\item \textbf{Paso 4 --- Filtro de Kalman.} Estima
      $\xv_t = (\text{presión}, \text{régimen})^\top$ a partir de
      $\zv_t = (\text{depth\_imb}, \text{ofi\_l1}, \text{vol\_30s})^\top$.
\item \textbf{Paso 5 --- Política Hedge.} Recibe $\xv_t$ y elige una
      acción muestreando de las probabilidades aprendidas (en la
      variante condicionada, las probabilidades dependen del bucket
      de presión).
\item \textbf{Paso 6 --- Simulador y pérdida.} Aplica la acción contra
      el libro actual y el siguiente, calcula slippage y adverse, y
      emite el escalar $L_t$ que actualiza los pesos del Hedge.
\end{enumerate}

%==============================================================================
\section{Pipeline de features}\label{sec:features}

El \texttt{FeaturePipeline} es \emph{stateful}: mantiene una
referencia al snapshot anterior (necesaria para OFI) y un buffer
circular de mids para la volatilidad (Sección~\ref{ssec:vol30s}).
Para cada snapshot del libro $\mathcal{B}_t$ produce un
\texttt{FeatureVector} con todos los campos descritos en la
Sección~\ref{sec:mercado}. La pipeline está diseñada para no
materializar el estado del libro en disco: los snapshots fluyen
directamente del \texttt{BookBuilder} a la pipeline. Solo el
\texttt{FeatureVector} se persiste en formato Parquet, particionado
por fecha y hora.

Cuando el filtro de Kalman recibe la observación, se quedan tres
campos del vector: $\text{depth\_imb}$, $\text{ofi\_l1}$,
$\text{vol\_30s}$. El resto del \texttt{FeatureVector}
($\text{spread\_abs}$, $\text{best\_bid}$, $\text{best\_ask}$,
$\text{mid}$, $\text{microprice}$) lo usa el simulador para calcular
slippage y adverse.

%==============================================================================
\section{Filtro de Kalman: deconstrucción matricial completa}\label{sec:kalman}

Esta sección desarma el filtro número por número. La meta es que el
lector pueda señalar cualquier celda de cualquier matriz y decir qué
hace.

\subsection{Modelo de espacio de estados}

El filtro asume un sistema lineal con estado oculto:
\begin{align}
\xv_t \;=\;& \Fm\,\xv_{t-1} + \mathbf{w}_t,\quad \mathbf{w}_t\sim\mathcal{N}(\mathbf{0},\Qm)
\label{eq:dynamics}\\[2pt]
\zv_t \;=\;& \Hm\,\xv_t + \mathbf{v}_t,\quad \mathbf{v}_t\sim\mathcal{N}(\mathbf{0},\Rm)
\label{eq:measurement}
\end{align}

\noindent
El estado $\xv_t \in \R^2$ y la observación $\zv_t \in \R^3$ son:
\begin{equation}
\xv_t = \begin{pmatrix}\text{pressure}_t \\ \text{regime}_t\end{pmatrix},
\qquad
\zv_t = \begin{pmatrix}\text{depth\_imb}_t \\ \text{ofi\_l1}_t \\ \text{vol\_30s}_t\end{pmatrix}.
\label{eq:state-obs}
\end{equation}

\subsection{Estandarización previa de la observación}

Antes de entrar al filtro, cada componente de $\zv_t$ se transforma como
\begin{equation}
\tilde{z}_{t,i} \;=\; \bigl(z_{t,i} - c_i\bigr)\, s_i,
\label{eq:standardize}
\end{equation}
con
\begin{equation}
\mathbf{c} = \begin{pmatrix}0\\0\\3{,}1658\!\cdot\!10^{-5}\end{pmatrix},\quad
\mathbf{s} = \begin{pmatrix}1\\1\\\num{36876,637}\end{pmatrix}.
\label{eq:standardize-vals}
\end{equation}

\noindent
\textbf{Por qué cada celda.}
\begin{itemize}
\item $c_0=c_1=0$, $s_0=s_1=1$: \texttt{depth\_imb} ya vive en el
      intervalo $(-1,+1)$ y \texttt{ofi\_l1} tiene desviación
      estándar empírica $\approx 1{,}51$. Estandarizarlas no aporta y
      sí destruiría las cotas (intencionales) de los buckets de presión.
\item $c_2 = 3{,}1658\!\cdot\!10^{-5}$ es el promedio empírico de
      \texttt{vol\_30s} en la primera ventana limpia de 38\,973 ticks
      (Apr~20--21).
\item $s_2 = 36876{,}637 = 1/\sigma$ donde $\sigma = 2{,}7117\!\cdot\!10^{-5}$
      es la desviación estándar empírica de \texttt{vol\_30s} en esa
      misma ventana.
\end{itemize}

\noindent
\textbf{Por qué hace falta.} Sin (\ref{eq:standardize}), $z_{t,2}$ vale
en magnitud $\sim\!10^{-5}$. Como $\Hm[2,1]=1$, la dimensión
\texttt{regime} del estado oculto está \emph{acotada} por la magnitud
de $z_{t,2}$ por construcción matemática: por más que se ajuste
$\Rm[2,2]$, regime no puede tener un rango dinámico útil si la
observación es del orden de $10^{-5}$. Con la estandarización
$\tilde{z}_{t,2}$ pasa a tener media cero y varianza unitaria, y
regime puede moverse en un rango interpretable del orden de unidades.
La verificación empírica está en la Sección~\ref{sec:resultados}:
después del fix la desviación de regime pasó de $3{,}17\!\cdot\!10^{-5}$
a $0{,}3709$.

\subsection{Matriz de transición $\Fm$}

\begin{equation}
\Fm = \begin{pmatrix}0{,}90 & 0{,}00 \\ 0{,}00 & 0{,}95\end{pmatrix}.
\label{eq:Fmatrix}
\end{equation}

\noindent
\textbf{Lectura celda por celda.}
\begin{itemize}
\item $\Fm[0,0]=0{,}90$: la presión en $t$ ``recuerda'' el 90\% de la
      presión en $t-1$. Equivale a un decaimiento exponencial con tiempo
      característico
      $\tau_{\text{pres}}\!\approx\!1/(1-0{,}90)\!=\!10$ ticks
      ($\sim\!1$ segundo a 100~ms por tick).
\item $\Fm[1,1]=0{,}95$: el régimen es más \emph{pegajoso}. Decaimiento
      característico de unos 20 ticks ($\sim\!2$ segundos). Este
      apareamiento refleja una creencia previa: la volatilidad cambia
      más despacio que la presión direccional.
\item $\Fm[0,1]=\Fm[1,0]=0$: presión y régimen no se acoplan en la
      transición. El filtro permite que la observación los acople a
      través de $\Hm$, pero la dinámica en sí los trata como dos canales
      independientes.
\end{itemize}

\subsection{Matriz de observación $\Hm$}

\begin{equation}
\Hm =
\begin{pmatrix}
1 & 0 \\
1 & 0 \\
0 & 1
\end{pmatrix}.
\label{eq:Hmatrix}
\end{equation}

\noindent
\textbf{Lectura celda por celda.}
\begin{itemize}
\item Filas 0 y 1 (\texttt{depth\_imb} y \texttt{ofi\_l1}) cargan
      $1$ en la columna 0 y $0$ en la columna 1: \emph{ambas observan
      directamente la presión y nada del régimen}. Esto es coherente
      con su semántica económica: ambas miden directamente desbalance
      direccional.
\item Fila 2 (\texttt{vol\_30s}) carga $0$ en la presión y $1$ en el
      régimen: la volatilidad realizada \emph{solo} informa sobre el
      estado de turbulencia, no sobre la dirección.
\item Sin entradas cruzadas: el filtro asume que ningún feature lleva
      información mezclada de presión y régimen. Esto simplifica la
      identificabilidad pero podría violarse si añadimos un feature
      mixto (e.g.\ momentum direccional ponderado por vol).
\end{itemize}

\subsection{Covarianza del proceso $\Qm$}

\begin{equation}
\Qm = \begin{pmatrix}0{,}010 & 0 \\ 0 & 0{,}005\end{pmatrix}.
\label{eq:Qmatrix}
\end{equation}

\noindent
\textbf{Lectura celda por celda.}
$\Qm[0,0]\!=\!0{,}010$ es la varianza del ruido aleatorio que puede
empujar la presión entre dos ticks (después del decaimiento por $\Fm$).
$\Qm[1,1]\!=\!0{,}005$ es la mitad: el régimen se mueve aleatoriamente
con menos amplitud, congruente con su mayor persistencia
($\Fm[1,1]\!=\!0{,}95$). Las entradas fuera de diagonal son cero porque
los dos procesos se modelan como ruido independiente.

\subsection{Covarianza de la observación $\Rm$}

\begin{equation}
\Rm = \begin{pmatrix}0{,}4486 & 0 & 0 \\ 0 & 2{,}2878 & 0 \\ 0 & 0 & 1{,}0\end{pmatrix}.
\label{eq:Rmatrix}
\end{equation}

\noindent
\textbf{Lectura celda por celda.}
\begin{itemize}
\item $\Rm[0,0]=0{,}4486$: varianza empírica de \texttt{depth\_imb}
      sobre la ventana limpia de 38\,973 ticks. Le dice al filtro:
      ``\texttt{depth\_imb} es ruidosa con una desviación estándar de
      $\sqrt{0{,}4486}\approx 0{,}67$''.
\item $\Rm[1,1]=2{,}2878$: varianza empírica de \texttt{ofi\_l1}.
      Equivale a $\sigma\approx 1{,}51$.
\item $\Rm[2,2]=1{,}0$: \emph{por construcción}, porque la
      estandarización (\ref{eq:standardize}) hace que la observación
      $\tilde{z}_{t,2}$ tenga ya varianza unitaria. Si más adelante se
      cambia el factor $s_2$, este valor también debe re-medirse.
\item Entradas fuera de diagonal cero: se asume que los ruidos de
      observación son independientes entre features. En la práctica
      tienen correlaciones pequeñas pero no se modelan; el filtro
      absorbe la diferencia en las ganancias.
\end{itemize}

\subsection{Estado y covarianza iniciales}

\begin{equation}
\xv_0 = \begin{pmatrix}0\\0\end{pmatrix},\qquad
\Pm_0 = \begin{pmatrix}1 & 0 \\ 0 & 1\end{pmatrix}.
\label{eq:initial}
\end{equation}

\noindent
\textbf{Por qué.}
$\xv_0$ vale cero porque al arrancar la sesión no se asume sesgo
direccional ni estado de régimen. $\Pm_0$ vale la identidad: se
declara una incertidumbre amplia, lo que hace que las primeras
observaciones tengan ganancias de Kalman altas y el filtro converja
rápido.

\subsection{Ciclo predict--update}

En cada tick el filtro aplica:

\medskip
\noindent
\textbf{Predict:}
\begin{align}
\hat{\xv}_{t|t-1} &= \Fm\,\xv_{t-1} \label{eq:predict-x}\\
\hat{\Pm}_{t|t-1} &= \Fm\,\Pm_{t-1}\,\Fm^\top + \Qm \label{eq:predict-P}
\end{align}

\medskip
\noindent
\textbf{Update:}
\begin{align}
\mathbf{y}_t &= \tilde{\zv}_t - \Hm\,\hat{\xv}_{t|t-1} \label{eq:innovation}\\
\Sm_t &= \Hm\,\hat{\Pm}_{t|t-1}\,\Hm^\top + \Rm \label{eq:innov-cov}\\
\Km_t &= \hat{\Pm}_{t|t-1}\,\Hm^\top\,\Sm_t^{-1} \label{eq:kalman-gain}\\
\xv_t &= \hat{\xv}_{t|t-1} + \Km_t\,\mathbf{y}_t \label{eq:posterior-x}\\
\Pm_t &= (\Iv - \Km_t\,\Hm)\,\hat{\Pm}_{t|t-1} \label{eq:posterior-P}
\end{align}

\medskip
\noindent
\textbf{Lectura intuitiva.}
La ecuación~(\ref{eq:innovation}) calcula el residual de la
predicción: cuánto se equivocó el filtro al pronosticar la observación.
La (\ref{eq:innov-cov}) es la covarianza esperada de ese residual.
La (\ref{eq:kalman-gain}) es la \emph{ganancia de Kalman}: una matriz
$2\!\times\!3$ que dice cuánto del residual se aplica al estado.
Si la observación es muy ruidosa (R grande) y la predicción muy
confiable (P pequeño), la ganancia es chica y el filtro confía en su
propia predicción. Si la predicción es incierta (P grande) y la
observación poco ruidosa (R pequeño), la ganancia es grande y el
filtro deja que la observación lo corrija.

\subsection{Tamaños y dimensiones (resumen)}

\begin{center}
\begin{tabular}{lcr}
\toprule
Símbolo & Forma & Significado \\
\midrule
$\xv_t$        & $2\!\times\!1$ & estado oculto (presión, régimen)\\
$\zv_t$        & $3\!\times\!1$ & observación cruda \\
$\tilde{\zv}_t$& $3\!\times\!1$ & observación estandarizada \\
$\Fm$          & $2\!\times\!2$ & transición de estado \\
$\Hm$          & $3\!\times\!2$ & matriz de observación \\
$\Qm$          & $2\!\times\!2$ & ruido de proceso \\
$\Rm$          & $3\!\times\!3$ & ruido de observación \\
$\Pm_t$        & $2\!\times\!2$ & covarianza posterior \\
$\Km_t$        & $2\!\times\!3$ & ganancia de Kalman \\
\bottomrule
\end{tabular}
\end{center}

%==============================================================================
\section{Política Hedge: algoritmo y variantes}\label{sec:hedge}

El filtro entrega cada tick un estado $\xv_t$. La política decide qué
hacer. Se evaluaron cuatro variantes; todas comparten el mismo núcleo
algorítmico.

\subsection{Algoritmo Hedge clásico (\emph{marginal})}

Hedge mantiene una distribución de probabilidad $\bm{\pi}_t$ sobre
las tres acciones $\mathcal{A}=\{\textsc{Wait},\textsc{Pass},\textsc{Agg}\}$.
En cada tick:
\begin{enumerate}
\item Muestrea $a_t \sim \bm{\pi}_t$.
\item Observa la pérdida $L_t$ asociada \emph{únicamente} a $a_t$
      (\emph{bandit feedback}).
\item Actualiza el peso de $a_t$:
\begin{equation}
w_{t+1}(a_t) \;=\; w_t(a_t) \cdot \exp\bigl(-\eta\,L_t\bigr),
\label{eq:hedge-update}
\end{equation}
con $\eta=0{,}1$ (learning rate).
\item Renormaliza para que $\sum_a w_{t+1}(a) = 1$, lo cual define
      $\bm{\pi}_{t+1}$.
\end{enumerate}

Las acciones que reciben pérdidas altas pierden peso exponencialmente;
las acciones con pérdidas bajas o negativas (ganancias) lo conservan
o crecen. La cota clásica de regret de Hedge se cumple por bucket:
$\sum_t L_t(a_t) - \min_a \sum_t L_t(a) \leq \mathcal{O}(\sqrt{T\,\log K})$
con $T$ ticks y $K\!=\!3$ acciones.

\medskip

\noindent
\textbf{Pesos iniciales.} $w_0(\textsc{Wait})=w_0(\textsc{Pass})=w_0(\textsc{Agg})=1/3$.
Distribución uniforme: ningún sesgo previo.

\subsection{Hedge condicionado por presión (1D bucketed)}

Esta es la versión que ganó el experimento. Discretiza la presión
$\xv_t[0]$ en seis buckets usando los cortes
\[
\text{pressure\_edges} = (-0{,}5,\, -0{,}2,\, 0,\, 0{,}2,\, 0{,}5),
\]
generando los buckets:

\begin{center}\small
\begin{tabular}{cl}
\toprule
Bucket & Rango de presión \\
\midrule
0 & $(-\infty,\, -0{,}5\,]$ \\
1 & $(-0{,}5,\, -0{,}2\,]$ \\
2 & $(-0{,}2,\, 0\,]$ \\
3 & $(0,\, 0{,}2\,]$ \\
4 & $(0{,}2,\, 0{,}5\,]$ \\
5 & $(0{,}5,\, +\infty)$ \\
\bottomrule
\end{tabular}
\end{center}

\medskip

Hay seis vectores de pesos $\{\bm{\pi}^{(b)}_t\}_{b=0}^{5}$, uno por
bucket. Al recibir $\xv_t$, se calcula $b_t = \mathrm{bucket}(\xv_t[0])$
con búsqueda binaria; se muestrea $a_t \sim \bm{\pi}^{(b_t)}_t$; se
observa $L_t$; y \emph{solo} $\bm{\pi}^{(b_t)}_t$ se actualiza con
(\ref{eq:hedge-update}). Las garantías clásicas de Hedge se cumplen
\emph{dentro} de cada bucket porque cada uno es un proceso de
exponential-weights independiente.

\subsection{Hedge condicionado 2D (presión $\times$ régimen)}

Igual que la 1D pero con dos ejes. Se añaden los cortes de régimen
$\text{regime\_edges}=(-0{,}3,\,0{,}3)$, generando 3 buckets de régimen
(calmo / normal / elevado). El total es $6\!\times\!3=18$ celdas. Se
mantiene un $\bm{\pi}^{(p,r)}_t$ por celda.

\subsection{Hedge condicionado adversarial (presión $\times$ señal-rolling)}

La hipótesis fue que el problema con la 2D regime no era el
condicionamiento bidimensional sino que el segundo eje (régimen
del Kalman) estaba detectando la \emph{probabilidad de fill}, no el
\emph{riesgo adverso}. Se reemplazó el segundo eje por una de tres
señales rolling de 60 segundos:
\begin{align}
\text{vol\_delta}_t &= \text{vol\_30s}_t - \text{vol\_30s}_{t-60s}\\
\text{ofi\_window}_t &= \sum_{s\in(t-60s,\,t]} \text{ofi\_l1}_s\\
\text{spread\_delta}_t &= \text{spread\_abs}_t - \text{spread\_abs}_{t-60s}
\end{align}
Los cortes del segundo eje se eligen por cuantiles empíricos pooled
sobre las 3 sesiones disponibles al momento del experimento (ver
Tabla~\ref{tab:cuts-secondaria}).

\begin{table}[!htbp]
\caption{Cortes del eje secundario por señal (cuantiles p33/p66
         pooled sobre las 3 sesiones).}
\label{tab:cuts-secondaria}
\centering\small
\begin{tabular}{lll}
\toprule
Señal & Cortes & Buckets \\
\midrule
\texttt{vol\_delta}    & $(-1{,}1\!\cdot\!10^{-5},\,+9{,}0\!\cdot\!10^{-6})$ & 3 \\
\texttt{ofi\_window}   & $(-2{,}0,\,+1{,}9)$                                 & 3 \\
\texttt{spread\_delta} & $(0)$                                               & 2 \\
\bottomrule
\end{tabular}
\end{table}

\begin{algorithm}[!htbp]
\caption{Hedge bucketed (versión 1D, pseudo-código).}
\label{alg:hedge1d}
\begin{algorithmic}[1]
\Require Cortes de presión, $\eta$, semilla RNG
\Ensure  Para cada tick $t$, una acción $a_t$ y un peso actualizado
\State Inicializar $\bm{\pi}^{(b)}_0 = (1/3, 1/3, 1/3)$ para $b=0\ldots 5$
\For{cada tick $t$}
  \State $\xv_t \gets$ filtro Kalman$(\zv_t)$
  \State $b_t \gets \textsc{bucket}(\xv_t[0])$
  \State $a_t \sim \bm{\pi}^{(b_t)}_t$
  \State $L_t \gets$ simulador$(a_t,\;\text{libro}_t,\;\text{libro}_{t+1})$
  \State $w \gets \bm{\pi}^{(b_t)}_t[a_t] \cdot e^{-\eta L_t}$
  \State $\bm{\pi}^{(b_t)}_t[a_t] \gets w$
  \State Renormalizar $\bm{\pi}^{(b_t)}_t$ para sumar 1
\EndFor
\end{algorithmic}
\end{algorithm}

%==============================================================================
\section{Simulador de ejecución y función de pérdida}\label{sec:simulador}

\subsection{Modelo de fill}

El simulador es determinista dadas $\text{mid}_t$, $\text{best\_bid}_t$,
$\text{best\_ask}_t$ y $\text{mid}_{t+1}$. Para cada acción:
\begin{itemize}
\item \textsc{Wait}: nunca ejecuta. $p^{\text{fill}} = 0$, $\text{filled}=\text{false}$.
\item \textsc{Passive}: ejecuta \emph{si y solo si} $\text{mid}_{t+1} \leq \text{mid}_t$,
       en cuyo caso $p^{\text{fill}} = \text{best\_bid}_t$. La regla
       captura la intuición de que un orden pasivo solo se cruza si
       aparece flujo vendedor descendente.
\item \textsc{Aggressive}: siempre ejecuta. $p^{\text{fill}} = \text{best\_ask}_t$.
\end{itemize}

\subsection{Función de pérdida y elección de $\lambda$}

La pérdida por tick es la ec.~(\ref{eq:loss}). El parámetro $\lambda$
es la palanca clave: pondera el costo de oportunidad
($\text{adverse\_move}$) frente al costo en dólares
($\text{slippage}$). Sobre datos limpios (38\,972 ticks, abril 2026):
\[
|\overline{\text{slip}}|=0{,}0013,\quad
\overline{\text{adv}}=0{,}4750,\quad
\sigma_{\text{slip}}=0{,}0247,\quad
\sigma_{\text{adv}}=2{,}3022.
\]

\noindent
Razones de balance:
\begin{itemize}
\item $\lambda=0{,}0027 = |\overline{\text{slip}}|/\overline{\text{adv}}$
      iguala las contribuciones medias.
\item $\lambda=0{,}0107 = \sigma_{\text{slip}}/\sigma_{\text{adv}}$
      iguala las contribuciones de ruido.
\item $\lambda=0{,}1$ (elegido) está entre 0{,}011 y 0{,}5 (default
      original). En este punto, $\textsc{Wait}$ y $\textsc{Aggressive}$
      tienen pérdida esperada en equilibrio prácticamente igual
      (0{,}040 vs 0{,}044) mientras que \textsc{Passive} se penaliza
      por selección adversa. Es decir, $\lambda=0{,}1$ obliga al
      Hedge a usar realmente la información de estado para discriminar
      entre las dos buenas acciones, en lugar de quedarse atascado en
      la marginal trivial.
\end{itemize}

\subsection{Configuraciones aplicadas}

\begin{table}[!htbp]
\caption{Hiperparámetros del experimento.}
\label{tab:hiperparams}
\centering\small
\begin{tabular}{lll}
\toprule
Parámetro & Valor & Donde está \\
\midrule
$\eta$ (Hedge)        & 0{,}10  & \texttt{configs/policy.yaml} \\
$\lambda$ (loss)      & 0{,}10  & \texttt{configs/execution.yaml} \\
\texttt{pressure\_edges} & $[-0{,}5,-0{,}2,0,0{,}2,0{,}5]$ & \texttt{configs/policy.yaml} \\
\texttt{regime\_edges}   & $[-0{,}3,\,0{,}3]$ & \texttt{configs/policy.yaml} \\
$\Fm,\Qm,\Rm,\Hm$       & ec.\,(\ref{eq:Fmatrix}--\ref{eq:Rmatrix}) & \texttt{configs/kalman.yaml} \\
$\xv_0,\Pm_0$           & ec.\,(\ref{eq:initial}) & \texttt{configs/kalman.yaml} \\
Semillas usadas         & $\{0,1,\ldots,9\}$ & --- \\
\bottomrule
\end{tabular}
\end{table}

%==============================================================================
\section{Diseño experimental}\label{sec:experimentos}

\subsection{Datos}

Se ingestaron en total cuatro sesiones limpias (sin huecos en la
secuencia de identificadores de actualización del WebSocket) entre
abril y mayo de 2026, ver Tabla~\ref{tab:sesiones}. Cada sesión es un
flujo continuo de eventos de \texttt{depthUpdate} y \texttt{trade}
desde Binance, persistido en JSONL. La feature pipeline produce un
parquet con un \texttt{FeatureVector} por tick.

\begin{table}[!htbp]
\caption{Sesiones de evaluación.}
\label{tab:sesiones}
\centering\small
\begin{tabular}{cllrr}
\toprule
ID & Inicio (UTC) & Fin (UTC) & Horas & Ticks \\
\midrule
S1 & 2026-04-20 19:54 & 2026-04-21 06:43 & 10{,}83 & \num{38973} \\
S2 & 2026-04-28 03:12 & 2026-04-28 14:02 & 10{,}83 & \num{38977} \\
S3 & 2026-04-30 01:45 & 2026-04-30 19:25 & 17{,}66 & \num{63590} \\
S7 & 2026-04-30 23:00 & 2026-05-01 20:05 & 21{,}08 & \num{75885} \\
\bottomrule
\end{tabular}
\end{table}

S4--S6 son fragmentos generados por el patch de \emph{self-heal} del
ingesto y se excluyen de todo análisis (entre 73 y 7\,060 filas
cada uno).

\subsection{Protocolo A/B pareado}

Para cada sesión y cada par (modo, semilla):
\begin{enumerate}
\item Se reinicializa el filtro de Kalman desde $(\xv_0,\Pm_0)$.
\item Se reinicializa la política con la semilla.
\item Se replay-ea la sesión completa, tick por tick: filtro
      $\to$ política $\to$ simulador $\to$ pérdida $\to$ actualización.
\item Se acumulan las métricas: total\_loss, slippage total, adverse
      total, fill\_rate, mezcla de acciones (porcentaje WAIT/PASS/AGG).
\end{enumerate}

Se usan diez semillas $(0,\ldots,9)$. La comparación es \emph{pareada}
por semilla: la diferencia $L^{\text{rico}}_s - L^{\text{1D}}_s$ se
calcula con $s$ misma para ambos modos. Esto cancela el ruido de
muestreo del RNG y reduce mucho la varianza del test.

\subsection{Estadística}

Para cada par (modo rico, sesión) reportamos:
\begin{itemize}
\item Media y desviación estándar de la diferencia pareada
      $d_s = L^{\text{rico}}_s - L^{\text{1D}}_s$.
\item Estadístico $t$ pareado: $|t| = |\bar d|/(\sigma_d/\sqrt{n})$,
      con $n=10$.
\item Conteo de semillas en las que el modo rico vence al 1D.
\end{itemize}
La interpretación tradicional es $|t|<2$ ``ruido'',
$|t|\!\approx\!2{-}3$ ``en el límite'', $|t|>3$ ``probablemente real'',
$|t|>10$ ``sin duda''.

\subsection{Regla de las dos sesiones}

Una lección aprendida durante el desarrollo: ningún competidor se
considera ganador con \emph{una sola sesión} a $|t|\!\geq\!2$, porque
con 16 comparaciones independientes (4 sesiones $\times$ 4 modos)
se espera que aproximadamente uno cruce $|t|=2$ por puro azar. La
regla aplicada en la fase final de evaluación es: \emph{se requiere
ganancia con $|t|\!\geq\!2$ y mismo signo en al menos dos sesiones
independientes antes de promover una variante}.

%==============================================================================
\section{Resultados}\label{sec:resultados}

\subsection{Resumen numérico}

La Tabla~\ref{tab:master} resume las quince comparaciones pareadas
producidas por los experimentos. Las cinco figuras
(Fig.~\ref{fig:fig01}--\ref{fig:fig05}) están generadas a partir del
JSON \texttt{ab\_adverse\_3sessions\_10seed.json} mediante
\texttt{scripts/\_plots\_section25.py}.

\begin{table*}[!htbp]
\caption{Resumen de las comparaciones pareadas. Una entrada
         ``$x/10$'' en \emph{wins} indica el número de semillas en
         las que la columna del título derrotó al 1D bucketed.
         La columna \emph{$|t|$} es el estadístico $t$ pareado.}
\label{tab:master}
\centering\small
\begin{tabular}{llccc}
\toprule
Test & Sesión & Ticks & $|t|$ & wins/10 (1D) \\
\midrule
1D vs marginal       & S1 & \num{38972} & 64{,}15 & 10/10 \\
\midrule
1D vs 2D regime      & S2 & \num{38976} & 17{,}73 & 10/10 \\
1D vs 2D regime      & S3 & \num{63589} & 37{,}67 & 10/10 \\
\midrule
1D vs vol\_delta     & S1 & \num{38972} & 28{,}89 & 10/10 \\
1D vs vol\_delta     & S2 & \num{38976} & 34{,}73 & 10/10 \\
1D vs vol\_delta     & S3 & \num{63589} & 34{,}72 & 10/10 \\
1D vs vol\_delta     & S7 & \num{75885} & 25{,}05 & 10/10 \\
\midrule
1D vs spread\_delta  & S1 & \num{38972} & 12{,}65 & 10/10 \\
1D vs spread\_delta  & S2 & \num{38976} & 12{,}20 & 10/10 \\
1D vs spread\_delta  & S3 & \num{63589} & 28{,}48 & 10/10 \\
1D vs spread\_delta  & S7 & \num{75885} & 24{,}42 & 10/10 \\
\midrule
1D vs ofi\_window    & S1 & \num{38972} & \phantom{0}2{,}97 & \phantom{0}8/10 \\
1D vs ofi\_window    & S2 & \num{38976} & \phantom{0}5{,}45 & \phantom{0}9/10 \\
1D vs ofi\_window    & S3 & \num{63589} & \phantom{0}2{,}37 & \phantom{0}2/10$^\dagger$ \\
1D vs ofi\_window    & S7 & \num{75885} & \phantom{0}0{,}23 & \phantom{0}5/10 \\
\bottomrule
\multicolumn{5}{l}{\footnotesize $^\dagger$ \texttt{ofi\_window} ganó 8/10 en S3.}
\end{tabular}
\end{table*}

\subsection{Figura 1 --- Pérdida total por sesión y modo}

\begin{figure*}[!htbp]
\centering
\includegraphics[width=0.92\textwidth]{fig01_total_loss_per_session.png}
\caption{Pérdida total media por (sesión $\times$ modo), con barras
         de error de desviación estándar a través de 10 semillas
         ($\lambda=0{,}1$). En cada grupo, la barra gris (1D baseline)
         es la más baja en S1 y S2; en S3 y S7 está empatada
         visualmente con \texttt{ofi\_window} (verde). Las barras
         azul (\texttt{vol\_delta}) y roja (\texttt{spread\_delta})
         son sistemáticamente más altas en las cuatro sesiones.}
\label{fig:fig01}
\end{figure*}

\textbf{Lectura.} La pérdida total escala aproximadamente lineal con
el número de ticks por sesión: 38k ticks $\to\!\sim\!1\,500$ unidades
de pérdida; 64k ticks $\to\!\sim\!2\,500$; 76k ticks $\to\!\sim\!2\,870$.
Esto confirma que el efecto por tick es estable y que las diferencias
absolutas entre modos crecen proporcionalmente con la longitud de la
sesión.

\subsection{Figura 2 --- Diferencias pareadas con etiquetas $|t|$}

\begin{figure*}[!htbp]
\centering
\includegraphics[width=0.92\textwidth]{fig02_paired_diff_vs_1d.png}
\caption{Diferencias pareadas (modo $-$ 1D) en la pérdida total, con
         estadístico $|t|$ etiquetado sobre cada barra. Barra positiva
         = el modo pierde frente a 1D. Solo \texttt{ofi\_window}
         (verde) cruza por debajo de cero, y solo en S3
         ($|t|=2{,}37$). En S7 vuelve a ruido ($|t|=0{,}23$).}
\label{fig:fig02}
\end{figure*}

\textbf{Lectura.} Esta figura es la prueba visual del fracaso de
\texttt{vol\_delta} y \texttt{spread\_delta}: ambas pierden con
$|t|>12$ en \emph{las cuatro} sesiones. Para \texttt{ofi\_window}, el
punto crítico es la comparación S3 vs S7: la única sesión donde gana
(S3) lo hace en el borde de la significancia
($|t|=2{,}37$); cuando se le da más datos (S7, +19\% ticks), la
ventaja se desploma a $|t|=0{,}23$. Si el problema fuera tamaño
muestral, más datos deberían haber ayudado, no perjudicado.

\subsection{Figura 3 --- Scatter pareado por semilla}

\begin{figure*}[!htbp]
\centering
\includegraphics[width=0.92\textwidth]{fig03_per_seed_scatter_ofi_window_vs_1d.png}
\caption{Comparación pareada por semilla de \texttt{ofi\_window}
         contra 1D. Cuatro paneles, uno por sesión. Cada punto
         representa una semilla; la línea $y=x$ es la frontera de
         empate. Punto debajo de la línea $\Rightarrow$
         \texttt{ofi\_window} venció a 1D para esa semilla. Conteos:
         S1: 2/10. S2: 1/10. S3: 8/10. S7: 5/10.}
\label{fig:fig03}
\end{figure*}

\textbf{Lectura.} Total a través de las cuatro sesiones: 16 victorias
de \texttt{ofi\_window} en 40 comparaciones pareadas. Una política
verdaderamente dominante mostraría 10/10 en cada sesión, como hizo el
1D bucketed contra el marginal en S1 ($|t|\!=\!64$).

\subsection{Figura 4 --- Distribución de acciones por modo}

\begin{figure*}[!htbp]
\centering
\includegraphics[width=0.92\textwidth]{fig04_action_mix_per_session.png}
\caption{Mezcla media de acciones por (sesión $\times$ modo). En las
         cuatro sesiones cualquier modo ``rico'' (azul, verde, rojo)
         exhibe un \emph{aumento sistemático} en \textsc{Passive}
         (entre $+1{,}4$ y $+2{,}5$ puntos porcentuales) y la
         consiguiente reducción en \textsc{Wait} y/o \textsc{Aggressive}
         frente al 1D baseline (gris).}
\label{fig:fig04}
\end{figure*}

\textbf{Lectura.} Esta es la figura más importante para entender
\emph{por qué} las variantes ricas pierden. La dirección del cambio
es la misma sin importar la señal del segundo eje: las tres señales
(\texttt{vol\_delta}, \texttt{ofi\_window}, \texttt{spread\_delta})
empujan al agente hacia más \textsc{Passive}. Esto sugiere que las
tres están detectando un mismo subconjunto de momentos --- los de
\emph{fill fácil} --- y que el algoritmo Hedge, sin saberlo, está
re-pesando hacia la acción que parece localmente más barata pero que
acumula mayor adverse.

\subsection{Figura 5 --- Frontera Pareto slippage vs adverse}

\begin{figure*}[!htbp]
\centering
\includegraphics[width=0.78\textwidth]{fig05_slippage_vs_adverse.png}
\caption{Frontera Pareto bidimensional: total\_slippage en USDT (eje
         x) contra total\_adverse\_move en unidades crudas (eje y).
         Lower-left = mejor en ambos ejes. Un punto por (sesión,
         modo). Las cuatro sesiones forman tres clusters por longitud
         (10{,}8h, 17{,}7h, 21{,}1h). Dentro de cada cluster, los modos
         se ordenan consistentemente: 1D (gris) abajo-derecha (más
         slippage, menos adverse); \texttt{vol\_delta} (azul)
         arriba-izquierda (menos slippage, más adverse).}
\label{fig:fig05}
\end{figure*}

\textbf{Lectura.} La pendiente positiva de cada cluster es la
firma de la \emph{barrera de toxicidad}: cualquier modo que ahorra
slippage paga proporcionalmente más adverse. Empíricamente, en estos
datos $\Delta\text{slip}\!\approx\!\Delta\text{adverse}/250$.
A $\lambda\!=\!0{,}1$ el balance económico es $\Delta\text{slip}=
-0{,}1\!\times\!\Delta\text{adverse}$, es decir cada \$1 de slippage
ahorrado vale \$10 de adverse. Como en la realidad cada \$1 ahorrado
viene con \$25 de adverse equivalente, todo movimiento a lo largo
de la pendiente es perdedor neto. El único punto que escapa
parcialmente es \texttt{ofi\_window} en S3 (cae estrictamente debajo
y a la izquierda del 1D-S3). En las otras siete combinaciones
(2 modos x 4 sesiones más \texttt{ofi\_window} en S1/S2/S7), todos los
puntos caminan sobre la misma frontera.

\subsection{Comportamiento por bucket (1D, S3, semilla 0)}

Para visibilidad, la Tabla~\ref{tab:bucket-1d} muestra los pesos
finales por bucket de la política 1D bucketed sobre S3. Buckets de
los extremos (0 y 5) están sub-visitados y permanecen casi uniformes;
los buckets centrales convergen claramente.

\begin{table}[!htbp]
\caption{Pesos finales por bucket de presión, política 1D, S3,
         semilla 0.}
\label{tab:bucket-1d}
\centering\small
\begin{tabular}{ccrrrr}
\toprule
b & rango & visitas & WAIT & PASSIVE & AGGR \\
\midrule
0 & $(-\infty,\,-0{,}500]$ &     43 & 0{,}3740 & 0{,}3351 & 0{,}2909 \\
1 & $(-0{,}500,-0{,}200]$  & 16\,479 & 0{,}8870 & 0{,}0504 & 0{,}0626 \\
2 & $(-0{,}200,\,0]$        & 17\,212 & 0{,}6754 & 0{,}1384 & 0{,}1862 \\
3 & $(0,\,+0{,}200]$        & 16\,309 & 0{,}3025 & 0{,}1751 & 0{,}5224 \\
4 & $(+0{,}200,+0{,}500]$   & 13\,456 & 0{,}0823 & 0{,}1039 & 0{,}8138 \\
5 & $(+0{,}500,+\infty)$    &     90 & 0{,}3369 & 0{,}3080 & 0{,}3551 \\
\bottomrule
\end{tabular}
\end{table}

\textbf{Lectura.} El bucket 1 (presión vendedora moderada) ha
aprendido a esperar en el 88{,}7\% de los casos. El bucket 4
(presión compradora moderada) ha aprendido a cruzar agresivamente en
el 81{,}4\% de los casos. La política condicionada está usando
correctamente la presión: es la prueba empírica directa de que
``saber la presión es suficiente''.

%==============================================================================
\section{Discusión: la barrera de toxicidad}\label{sec:discusion}

Tres resultados aparentemente independientes convergen en un único
mecanismo:

\begin{enumerate}
\item \textbf{Toda variante rica aumenta \textsc{Passive}}
      ($+1{,}4$--$2{,}5$ pp; Fig.~\ref{fig:fig04}).
\item \textbf{Toda variante rica ahorra slippage y paga más adverse}
      (Fig.~\ref{fig:fig05}).
\item \textbf{Las tres señales adversariales (\texttt{vol\_delta},
      \texttt{ofi\_window}, \texttt{spread\_delta}) tienen mecanismos
      económicos distintos pero fallan idénticamente.}
\end{enumerate}

\noindent
La explicación unificadora es que \emph{los tres ejes secundarios son
en realidad detectores de fill fácil, no de riesgo adverso}.
Cuando la volatilidad de corto plazo aumenta, cuando el spread se
ensancha, o cuando hay flujo neto sostenido en una dirección,
\emph{aparecen oportunidades de cruzar a precio favorable} y el
update multiplicativo de Hedge re-pesa hacia \textsc{Passive} en la
celda correspondiente. Pero la microestructura de los fills pasivos
en BTC/USDT es \emph{tóxica}: el flujo que cruza un bid pasivo es,
en promedio, el flujo que tiene una razón fundamental para empujar el
precio en su contra. La razón:
\[
\frac{\text{adverse}_{\textsc{Pass}}}{\text{adverse}_{\textsc{Agg}}} \;\approx\; 2.
\]
Es decir, los fills pasivos son seguidos por movimientos
desfavorables del precio aproximadamente al doble que los fills
agresivos.

\medskip

\noindent
\textbf{Por qué 1D bucketed escapa a la trampa.} Con solo seis
buckets, la política nunca aprende patrones finos: solo aprende la
estructura global ``presión vendedora $\Rightarrow$ esperar; presión
compradora $\Rightarrow$ cruzar''. Esa estructura coincide con la
estructura óptima desde el punto de vista de toxicidad: cuando hay
presión vendedora real, la peor decisión es comprar pasivo (los
vendedores te están alimentando justo antes de la caída). Como la
política 1D no puede ver \emph{micro-momentos} de fill fácil, no se
deja arrastrar al cebo.

\medskip

\noindent
\textbf{Implicancia teórica.} Adicionar dimensiones de estado al
condicionamiento no es libre: cada dimensión adicional aumenta la
varianza de la estimación por celda (la cantidad de updates por
celda decrece como $1/N$), y si la información que aporta la
dimensión adicional está correlacionada con un sesgo sistemático
(aquí: ``oportunidad de fill''), el modelo va a sobreajustar a ese
sesgo en lugar de a la métrica objetivo.

\medskip

\noindent
\textbf{Implicancia práctica.} Para el problema BTC/USDT a 100 ms con
$\lambda=0{,}1$, agregar \emph{cualquier} segundo eje al
condicionamiento estado-acción no mejora la pérdida combinada. Tres
intentos independientes con tres semánticas diferentes apuntan al
mismo techo. La conclusión razonable es que \emph{el espacio de
estados está agotado} para la familia de políticas bucketed.

%==============================================================================
\section{Trabajos relacionados y contexto industrial}\label{sec:related}

Esta sección responde tres preguntas concretas: (i)~de dónde sale la
idea de un filtro de Kalman, (ii)~de dónde sale la idea del algoritmo
Hedge, y (iii)~si la combinación específica que usamos en este trabajo
es algo común en la industria financiera o si es un híbrido propio.
La respuesta corta es: ninguna de las dos piezas es nueva, las dos
tienen décadas de literatura y uso operativo, pero combinarlas
exactamente como aquí --- Kalman para inferencia de estado y Hedge
bandit para selección de acción discreta --- es menos estándar de lo
que parece. Llegamos a esa combinación por un camino de restricciones
prácticas que se explican abajo.

\subsection{El filtro de Kalman: del programa Apolo al trading
            cuantitativo}

El filtro de Kalman aparece en 1960 con un objetivo muy distinto al
nuestro: estimar la posición y velocidad de un cohete a partir de
mediciones ruidosas de inercia y radar~\cite{kalman1960}. Su primera
aplicación de gran escala fue el sistema de guiado del programa
Apolo en la NASA. Desde entonces es la herramienta estándar para
estimación de estado en cualquier sistema dinámico lineal con ruido
gaussiano: GPS, control de drones, robótica móvil, biomédica,
predicción meteorológica.

\medskip

\noindent
La adopción en finanzas llega en los años 80 a través de la
econometría de series de tiempo. Dos libros marcaron la transición:
Harvey~\cite{harvey1989} formalizó cómo expresar modelos econométricos
clásicos (suavizamientos exponenciales, ARIMA, modelos estructurales)
como sistemas de espacio de estados resolubles con Kalman; Hamilton
\cite{hamilton1994} popularizó el enfoque entre los economistas
empíricos. A partir de ahí el filtro entró en finanzas cuantitativas
para tareas concretas:
\begin{itemize}
\item \textbf{Volatilidad latente.} Modelos tipo Stochastic Volatility
      asumen que la volatilidad es un proceso oculto que se infiere
      de los retornos observados. Cuando el modelo se linealiza
      adecuadamente, el filtro de Kalman (o sus variantes extendidas)
      hace la inferencia.
\item \textbf{Modelos de factores con cargas variables en el tiempo.}
      Si los $\beta$ de un modelo CAPM o multifactor cambian con el
      régimen, se modelan como estados ocultos y el filtro los
      actualiza tick a tick.
\item \textbf{Cointegración dinámica.} Para pares de activos que se
      mueven juntos pero con un spread que deriva, el spread se
      modela como estado oculto y se filtra.
\item \textbf{Microestructura de mercado.} Trabajos como
      \cite{cont2014} y, más cerca de nuestro problema,
      \cite{stoikov2018} sobre el microprecio, usan filtros de
      espacio-estado para inferir el ``precio justo'' subyacente al
      libro de órdenes a partir de las cotizaciones ruidosas.
\end{itemize}

\medskip

\noindent
\textbf{En la industria.} Es difícil saber qué corre exactamente en
los servidores de Renaissance Technologies, Two Sigma, Citadel o
Jane Street, porque ninguna firma publica su stack. Pero los
\emph{papers} y conferencias públicas de sus equipos de research
(NeurIPS, KDD, INFORMS) confirman que el espacio de estados
gaussiano y los filtros de Kalman son herramientas estándar para:
estimación de alfas latentes, descomposición de señal/ruido, y
módulos de predicción dentro de pipelines más grandes. Lo importante
es entender que no es una técnica exótica: es lo que un ingeniero
financiero entrena en su primer año de maestría.

\medskip

\noindent
\textbf{¿Por qué nosotros lo elegimos?} Tres razones, todas
prácticas.
\begin{enumerate}
\item \emph{Forma cerrada}. Cada paso del filtro es una multiplicación
      de matrices y una inversa de $3\!\times\!3$. No hay
      entrenamiento, no hay hiperparámetros aprendidos, no hay riesgo
      de no-convergencia. Las matrices vienen del dominio del problema.
\item \emph{Datos limitados}. La primera sesión limpia eran solo 38\,973
      ticks. Cualquier cosa con muchos parámetros entrenables
      (redes neuronales, modelos de Markov ocultos con K estados
      grandes) habría sufrido sobreajuste serio. Kalman tiene
      esencialmente un parámetro ajustable por feature ($\Rm[i,i]$),
      lo cual cabe cómodamente.
\item \emph{Diagnóstico}. Cuando algo se rompe, podemos imprimir la
      ganancia de Kalman, la covarianza posterior, y la innovación.
      Cada número tiene una interpretación. Esto fue clave para
      darnos cuenta, en sanity~§16, de que la dimensión de régimen
      estaba muerta porque la observación tenía magnitud $\sim\!10^{-5}$.
\end{enumerate}

\subsection{El algoritmo Hedge: de teoría de juegos a ejecución
            algorítmica}

Hedge tiene un linaje doble: por un lado viene del aprendizaje en
línea (online learning) y por otro de la teoría de juegos.
La idea de los \emph{multiplicative weights updates} ya estaba en
Littlestone--Warmuth~\cite{littlestone1994} (algoritmo de la
mayoría ponderada, 1989). Freund y Schapire en
1997~\cite{freund1997} formalizaron Hedge en el contexto de
boosting y demostraron la cota de regret
$\mathcal{O}(\sqrt{T\log K})$ que usamos en este trabajo. Auer
et~al.\ en 2002~\cite{auer2002} extendieron el resultado al
escenario \emph{bandit} --- en el que solo se observa la pérdida
de la acción tomada, no de todas, que es exactamente nuestro caso.
La síntesis canónica de toda la teoría está en
Cesa-Bianchi y Lugosi~\cite{cesabianchi2006}.

\medskip

\noindent
La aplicación a finanzas también es vieja. El trabajo seminal es
\emph{Universal Portfolios} de Cover~\cite{cover1991}, que en 1991
demostró que un portafolio actualizado con multiplicative weights
sobre acciones individuales tiene desempeño asintóticamente
comparable al mejor portafolio fijo en hindsight. Esto es
literalmente Hedge aplicado a selección de cartera. La línea
continúa con todo el campo de \emph{Online Portfolio Selection},
y ya en los 2000 hay un cambio: el problema deja de ser ``qué
cartera mantener'' y pasa a ser ``cómo ejecutar una orden grande''.

\medskip

\noindent
\textbf{Ejecución óptima como problema de control.} El estándar
clásico es Almgren--Chriss~\cite{almgren2001}, que formula la
ejecución de una orden grande como un problema de control óptimo
estocástico cerrado: dado un trade-off entre impacto de mercado y
riesgo de varianza temporal, la solución es una trayectoria de
ejecución de forma cerrada (a menudo cercana a TWAP/VWAP).
Obizhaeva--Wang~\cite{obizhaeva2013} extendió el modelo con
dinámicas de oferta y demanda. El libro de Cartea, Jaimungal y
Penalva~\cite{cartea2015} es la referencia pedagógica
moderna y cubre tanto modelos de control óptimo como filtros de
estado para ejecución algorítmica.

\medskip

\noindent
\textbf{Ejecución óptima como problema de aprendizaje.} La rama
``aprendamos directamente de los datos'' arranca con
Nevmyvaka--Feng--Kearns en 2006~\cite{nevmyvaka2006}: aplicaron
Q-learning a la ejecución óptima de órdenes en datos reales de
Nasdaq. De ahí en adelante, todo el reinforcement learning aplicado
a microestructura cae en este linaje. Spooner et~al.~\cite{spooner2018}
usaron RL para market making; Schnaubelt~\cite{schnaubelt2022} aplicó
deep RL específicamente al placement de órdenes límite en
criptomonedas. Hedge bandit (\emph{exp3}) es el primo más simple y
más interpretable de toda esa familia: en lugar de una red neuronal
con miles de parámetros aprende explícitamente una distribución
sobre acciones que evoluciona con la pérdida observada.

\medskip

\noindent
\textbf{En la industria.} Algoritmos de \emph{smart order routing}
en brokers (Goldman Sachs SIGMA-X, Morgan Stanley Trading Tools,
Citadel Securities) usan rutinas de online learning para decidir
en cada milisegundo a qué venue (NYSE, Nasdaq, dark pools) enviar
cada slice de una orden. La capa de decisión es típicamente
exp3-style con priors entrenados offline; la capa de ejecución
es Almgren--Chriss-like. En crypto, market makers como Wintermute,
Jump Trading, GTS, y los desks internos de los grandes exchanges
usan stacks similares, aunque proprietarios. El elemento común es
que \emph{la decisión discreta de ``qué hacer ahora''} se reduce a
un problema de bandit con feedback parcial, y la herramienta
estándar para eso son las familias de algoritmos derivadas de
Hedge / EXP3 / UCB.

\medskip

\noindent
\textbf{¿Por qué nosotros lo elegimos?} Cuatro razones, otra vez todas
prácticas.
\begin{enumerate}
\item \emph{Cota de regret cerrada}. La garantía
      $\sum_t L_t(a_t) - \min_a \sum_t L_t(a) \leq \mathcal{O}(\sqrt{T\log K})$
      vale sin asumir nada sobre la distribución de las pérdidas.
      Para datos reales que pueden ser no-estacionarios, eso importa.
\item \emph{No requiere modelo del entorno}. Los enfoques de
      control óptimo necesitan especificar dinámicas del mercado;
      las redes neuronales requieren grandes cantidades de data
      similar a la de entrenamiento. Hedge solo necesita una
      pérdida escalar por tick.
\item \emph{Bandit feedback es nativo}. En ejecución solo
      observamos qué pasó con la acción que tomamos; no sabemos
      qué habría pasado si hubiéramos esperado o cruzado en su
      lugar. Eso es exactamente la condición de
      bandit~\cite{auer2002}, lo que descarta de entrada
      cualquier algoritmo que asume full-information feedback.
\item \emph{Interpretabilidad}. Cada bucket de presión exhibe sus
      pesos finales en un vector de tres números. Si el bucket de
      presión vendedora moderada termina con
      $w(\textsc{Wait})=0{,}89$, sabemos exactamente qué aprendió y
      podemos auditarlo. Una red neuronal con la misma información
      es opaca.
\end{enumerate}

\subsection{La combinación Kalman + Hedge en la literatura}

La combinación específica de este trabajo --- un filtro de Kalman que
estima un estado oculto continuo, seguido por un Hedge bandit que
selecciona acción condicionado en una discretización de ese estado
--- no es una arquitectura estándar con un nombre propio en la
literatura. Sin embargo, sus componentes y la lógica de separar
``percepción'' (Kalman) de ``decisión'' (Hedge) tienen muchos
antecedentes:

\begin{itemize}
\item Cartea--Jaimungal~\cite{cartea2015} construyen pipelines del
      mismo estilo: filtro de espacio-estado (a veces Kalman, a veces
      filtro de partículas) para inferir intensidades latentes de
      flujo de órdenes, y luego una capa de control que toma
      decisiones de ejecución.
\item Avellaneda--Stoikov~\cite{avellaneda2008} ofrecen un esquema
      de market making donde la cotización es función explícita de
      la inferencia de inventario y volatilidad --- nuevamente,
      separación entre estimación y acción.
\item En la rama deep RL, los métodos \emph{actor-critic} con
      \emph{auxiliary state estimation tasks} (e.g.\ predicción de
      retornos a corto plazo como pérdida auxiliar para el critic)
      son una forma moderna de la misma separación.
\end{itemize}

\noindent
La diferencia con las prácticas industriales modernas es que la
mayoría usa redes neuronales profundas en alguno (o ambos) de los dos
bloques. Nosotros nos quedamos con la versión clásica por
tres motivos: (a)~tenemos pocos datos por sesión, (b)~queríamos cada
bloque debugueable de forma independiente, (c)~el problema es
suficientemente bien-condicionado como para no necesitar la
expresividad extra.

\subsection{El camino que nos llevó a esta arquitectura}

El recorrido conceptual del proyecto fue, en orden:

\begin{enumerate}
\item \emph{Pregunta inicial}: ¿se puede usar control adaptativo
      sobre datos reales de criptomercado? Esta pregunta es
      ingeniera, no financiera. La motivación viene de los cursos
      de control y de la observación de que el libro de órdenes
      tiene una dinámica de oferta/demanda observable y mensurable.
\item \emph{Restricción de datos}: lo que podemos capturar
      personalmente desde la API pública de Binance son ráfagas de
      10--20 horas de WebSocket más snapshots. Eso son entre 38\,000
      y 80\,000 ticks por sesión limpia. Cualquier algoritmo con
      más parámetros que datos por bucket queda descartado por
      sobreajuste.
\item \emph{Restricción de feedback}: en simulación de ejecución
      uno solo conoce el resultado de la acción que tomó. No hay
      ground-truth contrafactual. Eso descarta supervised learning
      directo y obliga a algo del tipo bandit.
\item \emph{Restricción de diagnóstico}: como es un proyecto de
      tesis, queremos poder señalar exactamente qué falla cuando
      falla. Las cajas negras son antipedagógicas. Eso favorece
      Kalman lineal y Hedge multiplicativo, ambos de forma cerrada.
\item \emph{Conclusión arquitectónica}: separar la inferencia
      del estado (Kalman) de la decisión (Hedge bucketed sobre el
      estado). Cada bloque se valida por separado.
      El Kalman se valida con sanity checks sobre la magnitud del
      estado, la innovación y la covarianza; el Hedge se valida
      con curvas de convergencia de pesos y comparaciones A/B.
\end{enumerate}

\noindent
Visto desde afuera, llegamos a una arquitectura que está en la
intersección de lo \emph{clásico de control} (Kalman) y lo
\emph{clásico de aprendizaje en línea} (Hedge). Eso es por
elección consciente: el campo se ha movido hacia deep RL en los
últimos diez años, pero para datos limitados y necesidad de
diagnóstico esa elección no es óptima. Nuestra contribución no es
inventar nada en el plano teórico --- ambas piezas tienen décadas
en sus respectivas literaturas --- sino mostrar empíricamente
qué funciona y qué no funciona cuando se las acopla sobre datos
microestructurales reales de BTC/USDT, y aislar el mecanismo
(la barrera de toxicidad de la Sección~\ref{sec:discusion}) que
impone un techo a esta familia de políticas.

%==============================================================================
\section{Conclusiones y trabajo futuro}\label{sec:conclusiones}

\subsection{Conclusiones}

\begin{enumerate}
\item La arquitectura Kalman--Hedge funciona end-to-end sobre datos
      reales de BTC/USDT a 100 ms. La integración pasa todos los
      sanity checks empíricos (signos correctos de slippage,
      convergencia de pesos, distribuciones de acciones razonables).
\item El condicionamiento sobre la presión (1D bucketed) es una
      mejora robusta y grande sobre la política sin estado:
      $-29\%$ en pérdida combinada en S1 con $|t|=64{,}15$ y 10/10
      semillas.
\item Tres intentos independientes de añadir un eje secundario
      (régimen de Kalman, \texttt{vol\_delta}, \texttt{spread\_delta},
      \texttt{ofi\_window}) fallaron en perder a 1D en al menos
      tres de cuatro sesiones. Solo \texttt{ofi\_window} ganó en
      una sesión (S3, $|t|=2{,}37$) y empató en otra (S7,
      $|t|=0{,}23$); por la regla de las dos sesiones no se promueve.
\item El mecanismo común del fallo es la \emph{barrera de toxicidad}:
      cualquier eje que detecta oportunidades de fill empuja al
      Hedge hacia más \textsc{Passive}, y los fills pasivos son
      sistemáticamente más tóxicos que los agresivos en este mercado.
\end{enumerate}

\subsection{Trabajo futuro}

La dirección recomendada para el siguiente capítulo de la tesis es
\textbf{condicionamiento por tamaño de orden}, que cambia el
\emph{espacio de acciones} en lugar del espacio de estado:
\begin{align}
\mathcal{A}_{\text{nuevo}} \;=\; \{\, &\textsc{Wait},\notag\\
                                   &\textsc{Pass}_{0{,}5},\textsc{Pass}_{1},\textsc{Pass}_{2},\notag\\
                                   &\textsc{Agg}_{0{,}5},\textsc{Agg}_{1},\textsc{Agg}_{2}\,\}.
\end{align}
La hipótesis es que dotar al agente de la capacidad de modular el
\emph{tamaño} de un fill pasivo (más pequeño en celdas neutras donde
acecha la toxicidad, más grande en celdas direccionales donde la
señal es real) crea una \emph{nueva frontera} en el plano slippage
vs adverse, en lugar de un nuevo punto sobre la frontera vieja.
La implementación es ligera: extender el enum de \texttt{Action} de 3
a 7 valores, escalar slippage y adverse proporcionalmente al tamaño,
y mantener Hedge clásico ($K\!=\!7$) sobre la nueva acción set, todo
condicionado en la presión 1D ya validada. La cota de regret cambia
solo en el factor $\sqrt{\log K}$, lo cual es despreciable
asintóticamente.

\medskip

\noindent
\textbf{Pasos secundarios.}
\begin{itemize}
\item Confirmar la victoria 1D vs marginal sobre S2, S3, S7
      (actualmente solo medida en S1). Tres corridas más a 10 semillas.
\item Probar las versiones \emph{invertidas} de \texttt{vol\_delta}
      y \texttt{spread\_delta}: cuando esas señales aumentan, ser
      \emph{más agresivo} en lugar de más pasivo. Las señales son
      claramente informativas ($|t|>12$ en magnitud); el signo de la
      reacción podría estar al revés.
\item Si el condicionamiento por tamaño tampoco rompe la barrera, el
      resultado negativo en sí es publicable: cuatro intentos
      mecanísticamente convergentes prueban que \emph{en BTC/USDT a
      100 ms y $\lambda=0{,}1$, no hay condicionamiento de estado
      en la familia bucketed que escape a la frontera Pareto}.
\end{itemize}

%==============================================================================
\section*{Reproducibilidad}

Toda la cadena de procesamiento está versionada en el repositorio
\texttt{adversity-is-controllable}. Los datos crudos se persisten en
\texttt{data/raw/} (JSONL) y los derivados en \texttt{data/derived/}
(Parquet). Las cinco figuras se regeneran con
\begin{verbatim}
PYTHONPATH=. python scripts/_plots_section25.py
\end{verbatim}
desde \texttt{data/derived/ab\_adverse\_3sessions\_10seed.json}.
La harness experimental es
\texttt{scripts/\_ab\_adverse\_3sessions.py}; ejecutarla replica las
160 corridas (4 sesiones $\times$ 10 semillas $\times$ 4 modos) en
$\sim\!200$~s sobre un CPU típico de laptop.

%==============================================================================
\section*{Agradecimientos}

A todos los que mantuvieron despiertas las máquinas de ingest las
noches del 20 de abril al 1 de mayo de 2026, y a los seis bugs de
libro corrupto que enseñaron a fondo el protocolo de sincronización
de Binance.

%==============================================================================
\begin{thebibliography}{9}

\bibitem{kalman1960}
R.~E.~Kalman,
``A New Approach to Linear Filtering and Prediction Problems,''
\emph{Trans.\ ASME, J.\ Basic Eng.}, vol.~82, no.~D, pp.~35--45, 1960.

\bibitem{freund1997}
Y.~Freund and R.~E.~Schapire,
``A decision-theoretic generalization of on-line learning and an application to boosting,''
\emph{J.\ Comput.\ Syst.\ Sci.}, vol.~55, no.~1, pp.~119--139, 1997.

\bibitem{cesabianchi2006}
N.~Cesa-Bianchi and G.~Lugosi,
\emph{Prediction, Learning, and Games}.
Cambridge, U.K.: Cambridge Univ.\ Press, 2006.

\bibitem{auer2002}
P.~Auer, N.~Cesa-Bianchi, Y.~Freund, and R.~E.~Schapire,
``The non-stochastic multi-armed bandit problem,''
\emph{SIAM J.\ Comput.}, vol.~32, no.~1, pp.~48--77, 2002.

\bibitem{bouchaud2018trades}
J.-P.~Bouchaud, J.~Bonart, J.~Donier, and M.~Gould,
\emph{Trades, Quotes and Prices: Financial Markets under the Microscope}.
Cambridge, U.K.: Cambridge Univ.\ Press, 2018.

\bibitem{cont2014}
R.~Cont, A.~Kukanov, and S.~Stoikov,
``The price impact of order book events,''
\emph{J.\ Financ.\ Econom.}, vol.~12, no.~1, pp.~47--88, 2014.

\bibitem{almgren2001}
R.~Almgren and N.~Chriss,
``Optimal execution of portfolio transactions,''
\emph{J.\ Risk}, vol.~3, pp.~5--40, 2001.

\bibitem{obizhaeva2013}
A.~A.~Obizhaeva and J.~Wang,
``Optimal trading strategy and supply/demand dynamics,''
\emph{J.\ Financ.\ Markets}, vol.~16, no.~1, pp.~1--32, 2013.

\bibitem{stoikov2018}
S.~Stoikov,
``The micro-price: a high-frequency estimator of future prices,''
\emph{Quant.\ Finance}, vol.~18, no.~12, pp.~1959--1966, 2018.

\bibitem{binancedocs}
Binance, ``Spot WebSocket API documentation,''
[Online]. Available: https://binance-docs.github.io/apidocs/spot/en/

\bibitem{harvey1989}
A.~C.~Harvey,
\emph{Forecasting, Structural Time Series Models and the Kalman Filter}.
Cambridge, U.K.: Cambridge Univ.\ Press, 1989.

\bibitem{hamilton1994}
J.~D.~Hamilton,
\emph{Time Series Analysis}.
Princeton, NJ: Princeton Univ.\ Press, 1994.

\bibitem{cover1991}
T.~M.~Cover,
``Universal portfolios,''
\emph{Math.\ Finance}, vol.~1, no.~1, pp.~1--29, 1991.

\bibitem{littlestone1994}
N.~Littlestone and M.~K.~Warmuth,
``The weighted majority algorithm,''
\emph{Inf.\ Comput.}, vol.~108, no.~2, pp.~212--261, 1994.

\bibitem{cartea2015}
\'A.~Cartea, S.~Jaimungal, and J.~Penalva,
\emph{Algorithmic and High-Frequency Trading}.
Cambridge, U.K.: Cambridge Univ.\ Press, 2015.

\bibitem{nevmyvaka2006}
Y.~Nevmyvaka, Y.~Feng, and M.~Kearns,
``Reinforcement learning for optimized trade execution,''
in \emph{Proc.\ 23rd Int.\ Conf.\ Mach.\ Learn.\ (ICML)}, 2006, pp.~673--680.

\bibitem{spooner2018}
T.~Spooner, J.~Fearnley, R.~Savani, and A.~Koukorinis,
``Market making via reinforcement learning,''
in \emph{Proc.\ 17th Int.\ Conf.\ Auton.\ Agents and Multi-Agent Syst.\ (AAMAS)}, 2018, pp.~434--442.

\bibitem{schnaubelt2022}
M.~Schnaubelt,
``Deep reinforcement learning for the optimal placement of cryptocurrency limit orders,''
\emph{Eur.\ J.\ Oper.\ Res.}, vol.~296, no.~3, pp.~993--1006, 2022.

\bibitem{avellaneda2008}
M.~Avellaneda and S.~Stoikov,
``High-frequency trading in a limit order book,''
\emph{Quant.\ Finance}, vol.~8, no.~3, pp.~217--224, 2008.

\end{thebibliography}

\end{document}
